% Einheit 3 – Punkte und Werte (Nr. 22–33)
\setcounter{aufgabe}{21}
\einheitenkopf[e3][3 Punkte und Werte]{Einheit 3 von 5 · Punkte und Werte}
\zweigzeile{Hier lernst du, Funktionswerte zu berechnen, Punkte zu prüfen und Nullstellen und Achsenschnittpunkte zu bestimmen · neu in diesem Jahr · P10 oft · baut auf: lineare Funktion (Einheit 2), Gleichungen lösen, negative Zahlen}

\begin{aufgabe}{Ich kann einen $x$-Wert in die Gleichung einsetzen.\newline Setze ein und schreibe die Rechnung auf. Rechne nicht aus.}
\begin{teilezwei}
\tz $f(x) = 4x + 3$, \ $x = 2$ & \tz $f(x) = 2x - 5$, \ $x = 6$ \\
\tz $f(x) = -3x + 1$, \ $x = 5$ & \tz $f(x) = 5x + 2$, \ $x = -1$ \\
\tz $f(x) = -x + 7$, \ $x = -3$ & \\
\end{teilezwei}
\end{aufgabe}

\verfahren{Funktionswerte berechnen}

\begin{aufgabe}{Ich kann einen Funktionswert berechnen.\newline Berechne den Funktionswert.}
\begin{teilezwei}
\tz $f(x) = 3x + 4$: \quad \feld{f(2)} & \tz $f(x) = 5x - 3$: \quad \feld{f(4)} \\
\tz $f(x) = 2x + 6$: \quad \feld{f(5)} & \tz $f(x) = 6x - 10$: \quad \feld{f(2)} \\
\tz $f(x) = 3x + 5$: \quad \feld{f(-2)} & \tz $f(x) = 4x - 1$: \quad \feld{f(2{,}5)} \\
\end{teilezwei}
\begin{teile}
\teil $f(x) = -2x + 3$ \quad \wertetabelle{x}{f(x)}{-2,-1,0,1,2}
\teil $f(x) = \frac{3}{4}x - 2$: \ Berechne den Funktionswert an der Stelle $-8$. (P10 2017 OS) \quad \feld{f(-8)}
\end{teile}
\end{aufgabe}

\verfahren{Den $x$-Wert zu einem Funktionswert bestimmen}

\begin{aufgabe}{Ich kann den $x$-Wert zu einem Funktionswert berechnen.\newline Für welches $x$ hat $f$ den angegebenen Wert? Löse die Gleichung.}
\rechenplatz{3}
\begin{gleichungsraster}[2]
\gl{f(x) = 2x + 3;\ \ f(x) = 11} & \gl{f(x) = 3x - 1;\ \ f(x) = 14} \\
\gl{f(x) = 5x + 2;\ \ f(x) = 22} & \gl{f(x) = 4x + 2;\ \ f(x) = -6} \\
\gl{f(x) = -2x + 7;\ \ f(x) = 1} & \gl{f(x) = 0{,}5x + 4;\ \ f(x) = 6} \\
\end{gleichungsraster}
\end{aufgabe}

\verfahren{Punktprobe}

\begin{aufgabe}{Ich kann prüfen, ob ein Punkt auf einer Geraden liegt.\newline Liegt der Punkt auf der Geraden? Rechne und kreuze an.}
\punktprobenliste{3x-2/A(2|4), 3x-2/B(1|2), -x+5/C({-1}|6), 2x+3/D(0|{-3}), 0{,}5x+1/E(3|2{,}5), -2x+1/F({-2}|{-3})}
\begin{teile}
\teil $y = -\frac{2}{3}x + 4$, \quad $P({-6}|8)$ \quad (P10 2026 FOR) \quad \janein
\end{teile}
\end{aufgabe}

\verfahren{Nullstellen und Achsenschnittpunkte}

\begin{aufgabe}{Ich kann die Nullstelle einer Geraden berechnen.\newline Berechne die Nullstelle.}
\rechenplatz{3}
\begin{gleichungsraster}[2]
\gl{f(x) = 2x - 6} & \gl{f(x) = 3x - 12} \\
\gl{f(x) = 4x - 20} & \gl{f(x) = 5x + 10} \\
\end{gleichungsraster}
\end{aufgabe}

\begin{aufgabe}{Ich kann die Nullstelle einer Geraden berechnen – weiter.}
\begin{gleichungsraster}[2]
\gl{f(x) = -4x + 8} & \gl{f(x) = -0{,}5x + 3} \\
\gl{f(x) = 1{,}2x - 3 \quad \text{(P10 2022 OS)}} & \\
\end{gleichungsraster}
\end{aufgabe}

\begin{aufgabe}{Ich kann die Schnittpunkte einer Geraden mit den Achsen angeben.\newline Gib die Schnittpunkte mit der $y$-Achse und mit der $x$-Achse an.}
\begin{teile}
\teil $f(x) = 2x - 4$ \hfill $y$-Achse: \punktfeld \quad $x$-Achse: \punktfeld
\teil $f(x) = -3x + 6$ \hfill $y$-Achse: \punktfeld \quad $x$-Achse: \punktfeld
\teil $f(x) = 0{,}5x + 2$ \hfill $y$-Achse: \punktfeld \quad $x$-Achse: \punktfeld
\teil $f(x) = 4x$ \hfill $y$-Achse: \punktfeld \quad $x$-Achse: \punktfeld
\teil $f(x) = -2$ \hfill $y$-Achse: \punktfeld \quad $x$-Achse: \punktfeld
\end{teile}
\end{aufgabe}

\begin{aufgabe}{Ich kann Nullstellen und Schnittpunkte am Graphen ablesen.\newline Lies am Graphen ab.}

\begin{ksys}[xmin=-3,xmax=5,ymin=-4,ymax=5,ablesen]
\gerade[4.5]{1}{-1}{p}
\gerade[-1.5]{-1}{3}{q}
\end{ksys}

\begin{teilezwei}
\tz Nullstelle von $p$: \quad \feld{x} & \tz Nullstelle von $q$: \quad \feld{x} \\
\tz $q$ schneidet die $y$-Achse in \punktfeld & \tz $p$ und $q$ schneiden sich in \punktfeld \ (P10 2023 OS) \\
\end{teilezwei}
\end{aufgabe}

\verfahren{Prüfen, begründen, anwenden}

\begin{aufgabe}{Ich kann Aussagen über eine Gerade prüfen.\newline Die Gerade gehört zu $f(x) = -1{,}5x + 3$. Kreuze an: wahr oder falsch? (P10 2021 OS)}
\begin{teile}
\teil Die Gerade fällt. \hfill \kreuz{wahr} \kreuz{falsch}
\teil Die Gerade schneidet die $y$-Achse im Punkt $(0|{-1{,}5})$. \hfill \kreuz{wahr} \kreuz{falsch}
\teil Die Nullstelle ist $x = 2$. \hfill \kreuz{wahr} \kreuz{falsch}
\teil Der Punkt $(4|3)$ liegt auf der Geraden. \hfill \kreuz{wahr} \kreuz{falsch}
\end{teile}
\end{aufgabe}

\begin{aufgabe}{Ich finde den Fehler bei Nullstelle und Punktprobe.\newline Finde den Fehler, benenne ihn und korrigiere die Rechnung.}
\begin{teile}
\teil Jonas sucht die Nullstelle von $f(x) = 3x - 9$:
\rechnung{f(0) &= 3 \cdot 0 - 9 = -9 \\ \text{Nullstelle: } x &= -9}
\teil Jonas prüft, ob $P(5|2)$ auf dem Graphen von $f(x) = 2x - 8$ liegt:
\rechnung{f(2) &= 2 \cdot 2 - 8 = -4 \\ -4 &\neq 5 && \text{P liegt nicht auf dem Graphen.}}
\end{teile}
\end{aufgabe}

\begin{aufgabe}{Ich kann begründen, ob und wo eine Gerade eine Nullstelle hat.\newline Entscheide ohne Rechnung und begründe.}
\begin{teile}
\teil Hat $f(x) = 4$ eine Nullstelle?
\teil Liegt die Nullstelle von $f(x) = 2x + 6$ links oder rechts von der $y$-Achse?
\end{teile}
\end{aufgabe}

\begin{aufgabe}{Ich kann mit Funktionswert und Nullstelle eine Sachaufgabe lösen.\newline Ein Handyakku ist zu 90\,\% geladen. Beim Spielen sinkt die Ladung gleichmäßig. $y = -7{,}5x + 90$ gibt die Ladung in Prozent nach $x$ Stunden an.}
\begin{teile}
\teil Wie viel Prozent zeigt der Akku nach 2 Stunden? \quad \leerfeld[\%]
\teil Nach wie vielen Stunden zeigt er noch 30\,\%? \quad \leerfeld[h]
\teil Nach wie vielen Stunden ist der Akku leer? (P10 2021 OS) \quad \leerfeld[h]
\end{teile}
\end{aufgabe}
